\documentclass[aspectratio=169,10pt,spanish]{beamer}

\input{../../../_preambulo_beamer.tex}
\input{../../../_comandos_beamer.tex}

\selectlanguage{spanish}

\title{MA1001B - Modelación Estadística para la Toma de Decisiones}
\subtitle{Lección 47: Regresión lineal simple: mínimos cuadrados y $R^2$}
\author{}
\institute{Tecnol\'ogico de Monterrey}
\date{}

\begin{document}

\begin{frame}[plain]
  \vspace{1.2cm}
  {\Large\bfseries MA1001B - Modelación Estadística para la Toma de Decisiones\par}
  \vspace{0.7cm}
  {\large Lección 47: Regresión lineal simple: mínimos cuadrados y $R^2$\par}
  \vspace{0.7cm}
  {\color{TecRojo}\rule{\textwidth}{1.2pt}}
  \vspace{0.8cm}

  Facultad MA1001B\\[0.3cm]
  Periodo\\[0.3cm]
  Tecnol\'ogico de Monterrey
\end{frame}

\begin{frame}{La pregunta de ajustar una recta}
  \begin{alertblock}{Pregunta guía}
    ¿Cómo puede un predictor numérico resumirse con una recta única, y por
    qué una $R^2$ alta sigue sin ser una palanca causal de coaching?
  \end{alertblock}

  \vspace{0.6em}
  Al final de esta lección, usted puede:
  \begin{itemize}
    \item leer una nube de horas de coaching y unidades por hora;
    \item calcular la pendiente y el intercepto de mínimos cuadrados;
    \item formar $R^2=1-SSE/SST$;
    \item rechazar tratar una $R^2$ alta como un botón de política.
  \end{itemize}

  \vfill
  \scriptsize Todas las semanas y valores de productividad usados hoy son sintéticos. No se usan datos reales de empresa.
\end{frame}

\begin{frame}{Cuarenta semanas sintéticas}
  \small
  \begin{center}
    \begin{tabular}{lr}
      \toprule
      \textbf{Cantidad} & \textbf{Valor} \\
      \midrule
      Semanas $n$ & 40 \\
      Media de horas de coaching $\bar x$ & 25.075277 \\
      Media de unidades por hora $\bar y$ & 23.315500 \\
      Correlación muestral $r$ & 0.910813 \\
      Recta verdadera & $y=12+0.45x+\varepsilon$ \\
      \bottomrule
    \end{tabular}
  \end{center}

  \vspace{0.5em}
  \begin{exampleblock}{Muestra observacional}
    $x$ es horas de coaching esa semana. $y$ es unidades empacadas por
    hora-labor. El coaching se observó, no se asignó al azar.
  \end{exampleblock}
\end{frame}

\begin{frame}[fragile]{Paso 1: Nube, aún sin recta}
  \begin{columns}[T]
    \begin{column}{0.40\textwidth}
      \small
      \begin{lstlisting}
x = coaching_hours
y = units_per_hour
r = corrcoef(x, y)
      \end{lstlisting}

      \begin{block}{Salida verificada}
        $n=40$, $r=0.910813$\\
        $\bar x=25.075277$\\
        $\bar y=23.315500$
      \end{block}
    \end{column}
    \begin{column}{0.60\textwidth}
      \centering
      \includegraphics[width=\textwidth]{../figuras/regresion_simple_01_dispersion.png}
      \vspace{0.2em}
      \scriptsize Nube de 40 semanas del Paso 1
    \end{column}
  \end{columns}
\end{frame}

\begin{frame}{Los mínimos cuadrados minimizan huecos verticales}
  \small
  \[
    b_1=\frac{S_{xy}}{S_{xx}}=0.405936,
    \qquad
    b_0=\bar y-b_1\bar x=13.136554.
  \]
  Recta ajustada: $\hat y=13.136554+0.405936\,x$. Pendiente verdadera
  $0.45$; intercepto verdadero $12$. $SSE=113.102672$. scipy
  \texttt{linregress} coincide con la pendiente y el intercepto manuales.
\end{frame}

\begin{frame}[fragile]{Paso 2: La recta ajustada}
  \begin{columns}[T]
    \begin{column}{0.40\textwidth}
      \small
      \begin{lstlisting}
b1 = Sxy / Sxx
b0 = ybar - b1 * xbar
yhat = b0 + b1 * x
      \end{lstlisting}

      \begin{block}{Salida verificada}
        $b_1=0.405936$\\
        $b_0=13.136554$\\
        $SSE=113.102672$
      \end{block}
    \end{column}
    \begin{column}{0.60\textwidth}
      \centering
      \includegraphics[width=\textwidth]{../figuras/regresion_simple_02_minimos_cuadrados.png}
      \vspace{0.2em}
      \scriptsize Recta de mínimos cuadrados del Paso 2
    \end{column}
  \end{columns}
\end{frame}

\begin{frame}{Paso 3: Una $R^2$ alta no es una palanca causal}
  \begin{columns}[T]
    \begin{column}{0.42\textwidth}
      \small
      $SST=663.669939$, $SSR=550.567267$, $R^2=0.829580$.

      \vspace{0.4em}
      Ganancia ajustada de $+10$ horas: $4.059355$ unidades por hora.

      \vspace{0.5em}
      \textbf{Límite:} una $R^2$ alta no implica una palanca causal. El
      coaching no se aleatorizó. La Lección 48 añade $EE(b_1)$, una prueba
      $t$ y un punto influyente.
    \end{column}
    \begin{column}{0.58\textwidth}
      \centering
      \includegraphics[width=\textwidth]{../figuras/regresion_simple_03_limite_r_cuadrada.png}
      \vspace{0.2em}
      \scriptsize Descomposición de $R^2$ del Paso 3
    \end{column}
  \end{columns}
\end{frame}

\begin{frame}{Verificación de conceptos}
  \begin{enumerate}
    \item ¿Qué cantidad vertical minimizan los mínimos cuadrados?
    \item ¿Por qué $b_1=0.405936$ puede diferir de la pendiente verdadera
      $0.45$?
    \item Verifique $SSR+SSE=SST$ a partir de las sumas de cuadrados
      impresas.
    \item ¿Por qué $R^2=0.829580$ no es razón para mandar 10 horas más de
      coaching?
  \end{enumerate}

  \vspace{0.6em}
  \begin{block}{Conexión con la Lección 48}
    La sesión puede continuar directamente con inferencia de la pendiente,
    gráficas de residuos y puntos influyentes.
  \end{block}

  \vfill
  \scriptsize Fuente: OpenStax, \textit{Introductory Business Statistics 2e},
  Secciones 13.3--13.4. CC BY-NC-SA.
\end{frame}

\appendix



\end{document}
